\subsection{Recurrent user preference}
\begin{frame}{Stage 2 — Preference như một trạng thái recurrent}
  \begin{columns}[T,onlytextwidth]
    \begin{column}{0.56\textwidth}\centering
      \includegraphics[width=0.98\textwidth]{recurrent_preference.png}
    \end{column}
    \begin{column}{0.4\textwidth}\small
      Chia lịch sử thành các block có thứ tự:
      \begin{equation*}
        \begin{aligned}
          h_k&=\mathrm{LLM}(\mathrm{summary}(B_k),h_{k-1}),\\[-0.1cm]
          h_0&=\varnothing.
        \end{aligned}
      \end{equation*}
      \begin{itemize}
        \item Block đầu tạo preference ban đầu.
        \item Block sau đọc preference trước đó và cập nhật.
        \item $h_k$ là văn bản, nên dễ inspect hơn embedding ẩn.
      \end{itemize}
    \end{column}
  \end{columns}
\end{frame}

\begin{frame}{Một vòng cập nhật preference diễn ra thế nào?}
  \begin{enumerate}
    \item Lấy 3 item liên tiếp (paper chọn block size 3 trong implementation).
    \item Mỗi item đã có title + visual description + item summary.
    \item Prompt yêu cầu LLM chỉ tóm tắt sở thích, không lặp lại chi tiết item.
    \item Kết quả trở thành context cho block kế tiếp.
  \end{enumerate}
  \vspace{0.25cm}\begin{block}{Trade-off}
    Block nhỏ: cập nhật thường xuyên nhưng nhiều lần gọi model. Block lớn: ít call hơn nhưng prompt dài, dễ làm mờ sự thay đổi theo thời gian.
  \end{block}
\end{frame}
